% Diffraction d'une lumière laser
\begin{minipage}{0.42\linewidth}
    La lumière issue d'un laser est diffractée par une fente fine de largeur $a$. La
    figure de diffraction est observée sur un écran situé à la distance $D$ de la
    fente.
    
    La largeur de la tache centrale a pour valeur $d$.
    
    On appelle $\alpha$ l'écart angulaire entre les rayons diffractés extrêmes
    correspondant à la tache centrale.
\end{minipage}
\hfill
\begin{minipage}{0.52\linewidth}
\begin{figure}[H]
    \centering
    \includegraphics[width=\textwidth]{2025-12-19_235735-}
\end{figure}
\end{minipage}

\begin{enumerate}
    \item L'écran est situé loin de la fente et l'angle $\alpha$ est petit. On
          peut donc écrire que tan $\alpha\approx\alpha$ ($\alpha$ étant
          exprimé en radians). En déduire l'expression de $\alpha$ en fonction
          de $d$ et $D$.
    \item L'écart angulaire $\alpha$ varie en fonction de la longueur d'onde
          $\lambda$ de la radiation utilisée et de la largeur $a$ de la fente
          suivant la relation~: $\alpha=\frac{2\lambda}{a}$.

          La fente du dispositif décrit précédemment a pour largeur
          $\qty{0.10}{\mm}$.
    
          Elle est située à une distance $D=\qty{2.0}{\m}$ de l'écran, et on
          observe une tache centrale de largeur $d=\qty{2.1}{\cm}$.
          \begin{enumerate}
              \item Déterminer la longueur d'onde de la radiation
                    monochromatique émise par le laser.
              \item À quelle couleur correspond-elle~?
          \end{enumerate}
\end{enumerate}

